\chapter{Glossary}
\index{glossary}
\label{app:glossary}

\begin{description}[leftmargin=2em, style=nextline]

\item[alist]\index{alist}
An \emph{association list}: a list of pairs used as a key-value map. Keys are in the \texttt{car} of each pair, values in the \texttt{cdr}. Lookup is O(n) via \texttt{assoc} or \texttt{assq}.

\item[atom]\index{atom}
A value that is not a pair: numbers, strings, symbols, characters, booleans, vectors, and bytevectors are all atoms.

\item[bignum]\index{bignum}
An arbitrary-precision integer. Kaappi automatically promotes fixnums to bignums on overflow.

\item[bytecode]\index{bytecode}
The low-level instruction format that Kaappi's compiler produces. Executed by the register-based virtual machine.

\item[bytevector]\index{bytevector}
A fixed-length sequence of bytes (integers 0--255). Created with \texttt{make-bytevector} or \texttt{bytevector}. Used for binary data and I/O buffers.

\item[call/cc]\index{call/cc}
Short for \texttt{call-with-current-continuation}. Captures the current continuation as a first-class procedure that can be invoked later.

\item[car]\index{car}
The first element of a pair. From ``Contents of the Address Register'' on the IBM 704.

\item[cdr]\index{cdr}
The second element of a pair. From ``Contents of the Decrement Register.'' For a proper list, the cdr is the rest of the list.

\item[channel]\index{channel}
A synchronized communication primitive for fibers. Created with \texttt{make-channel}. \texttt{channel-send} and \texttt{channel-receive} block until both sender and receiver are ready.

\item[closure]\index{closure}
A function bundled with the environment in which it was defined. Closures capture free variables from their enclosing scope.

\item[cond-expand]\index{cond-expand}
A form for conditional compilation. Selects code branches based on features available in the implementation.

\item[condition]\index{condition}
A value raised as an exception. Can be an error object (created by \texttt{error}) or any Scheme value (raised by \texttt{raise}).

\item[cons cell]\index{cons cell}
A pair: a heap-allocated object with a \texttt{car} and a \texttt{cdr}. The fundamental building block of lists.

\item[continuation]\index{continuation}
The ``rest of the computation'' from a given point in the program. Scheme allows capturing continuations as first-class values.

\item[datum]\index{datum}
Any Scheme value that can be read by the reader: a number, string, symbol, list, vector, etc. A datum is the external representation of a value.

\item[define-library]\index{define-library}
The R7RS form for creating a library (module). Specifies exports, imports, and the library body.

\item[dotted pair]\index{dotted pair}
A pair whose cdr is not a list. Printed as \texttt{(a . b)}.

\item[dynamic-wind]\index{dynamic-wind}
A mechanism that guarantees setup and cleanup thunks are called when control enters or leaves a dynamic extent, even via continuations.

\item[eof object]\index{eof object}
A special value returned by input procedures when end-of-file is reached. Tested with \texttt{eof-object?}.

\item[exact]\index{exact}
A number whose value is represented precisely (integers, rationals). Contrasted with inexact (floating-point).

\item[export]\index{export}
A declaration in a library that makes a binding available to importers. Specified with the \texttt{export} clause in \texttt{define-library}.

\item[FFI (Foreign Function Interface)]\index{FFI}\index{Foreign Function Interface}
A mechanism for calling C functions from Scheme. Kaappi provides \texttt{(kaappi ffi)} with \texttt{ffi-open}, \texttt{ffi-fn}, and \texttt{ffi-callback}.

\item[fiber]\index{fiber}
A lightweight green thread managed by Kaappi's runtime. Created with \texttt{spawn} from \texttt{(kaappi fibers)}. Fibers are cooperatively scheduled.

\item[fixnum]\index{fixnum}
A small integer stored inline (no heap allocation) in Kaappi's NaN-boxed value representation. 63-bit range.

\item[flonum]\index{flonum}
An IEEE 754 double-precision floating-point number. Inexact by definition.

\item[form]\index{form}
A Scheme expression, particularly one with special syntax (e.g., \texttt{if}, \texttt{define}, \texttt{let}).

\item[guard clause]\index{guard clause}
An exception-handling form that catches raised exceptions and dispatches on conditions, similar to \texttt{try}/\texttt{except} in Python.

\item[hash table]\index{hash table}
A mutable key-value data structure with O(1) average lookup. Provided by SRFI-69. Created with \texttt{make-hash-table}.

\item[hygienic macro]\index{hygienic macro}
A macro system where generated identifiers cannot accidentally capture variables from the macro's call site. Scheme's \texttt{syntax-rules} is hygienic.

\item[identifier]\index{identifier}
A name in Scheme source code. Can contain letters, digits, and extended characters like \texttt{-}, \texttt{?}, and \texttt{!}.

\item[improper list]\index{improper list}
A chain of pairs not terminated by the empty list. The last cdr is a non-list value.

\item[inexact]\index{inexact}
A number whose value may be approximate (floating-point). Created by decimal literals or operations that introduce rounding.

\item[irritants]\index{irritants}
Additional values attached to an error object, providing context about what went wrong.

\item[lexical scope]\index{lexical scope}
A scoping rule where a name refers to the binding in the nearest enclosing scope where it was defined. Also called static scope.

\item[NaN-boxing]\index{NaN-boxing}
Kaappi's value representation technique: all Scheme values are stored in 64-bit words using the unused bit patterns of IEEE 754 NaN values.

\item[pair]\index{pair}
See \emph{cons cell}.

\item[port]\index{port}
An abstraction for I/O sources and sinks. Input ports read data; output ports write data. Ports can be connected to files, strings, or the console.

\item[predicate]\index{predicate}
A procedure that returns \texttt{\#t} or \texttt{\#f}. By convention, predicates end with \texttt{?} (e.g., \texttt{null?}, \texttt{string?}).

\item[procedure]\index{procedure}
A first-class callable value. Created by \texttt{lambda}, \texttt{define}, or built-in.

\item[promise]\index{promise}
A deferred computation created by \texttt{delay}. Its value is computed at most once, when \texttt{force} is called.

\item[proper list]\index{proper list}
A chain of pairs terminated by the empty list \texttt{'()}. The standard list type in Scheme.

\item[quasiquote]\index{quasiquote}
Template syntax using backtick (\texttt{`}) that allows selective unquoting with \texttt{,} and splicing with \texttt{,@}.

\item[R7RS]\index{R7RS}
The Revised$^7$ Report on the Algorithmic Language Scheme (2013). The current Scheme standard. ``R7RS-small'' is the core language; ``R7RS-large'' is the extended library effort.

\item[record]\index{record}
A user-defined structured data type with named fields. Created with \texttt{define-record-type}.

\item[S-expression]\index{S-expression}
The syntactic representation of Scheme code and data: atoms and nested lists in parentheses.

\item[sandbox]\index{sandbox}
A restricted execution mode (\texttt{--sandbox}) that disables file I/O, FFI, and network access for untrusted code.

\item[shared library]\index{shared library}
A compiled C library (\texttt{.dylib} on macOS, \texttt{.so} on Linux) that can be loaded at runtime via the FFI.

\item[spawn]\index{spawn}
The procedure that creates a new fiber from a thunk. The fiber runs concurrently with other fibers on the same OS thread.

\item[SRFI]\index{SRFI}
Scheme Requests for Implementation. Community-designed portable libraries. Kaappi ships 72 SRFIs (8 built-in, 64 portable).

\item[symbol]\index{symbol}
An interned identifier value. Two symbols with the same name are pointer-equal (\texttt{eq?}). Used as keys, tags, and in code representation.

\item[syntax-rules]\index{syntax-rules}
The pattern-based macro definition form in R7RS. Matches input against patterns and produces output from templates.

\item[tail call]\index{tail call}
A procedure call in tail position---the last expression evaluated before returning. Scheme guarantees tail calls use constant stack space.

\item[tail position]\index{tail position}
A syntactic position where a procedure call is the last action before returning. Calls in tail position are guaranteed not to grow the stack (tail-call optimization).

\item[TCO]\index{TCO}
Tail Call Optimization. The guarantee that tail calls do not consume stack space.

\item[thottam]\index{thottam}
Kaappi's built-in package manager. Installs, manages, and resolves dependencies for Kaappi ecosystem libraries.

\item[thunk]\index{thunk}
A zero-argument procedure, used to delay evaluation. \texttt{(lambda () expr)}.

\item[variadic]\index{variadic}
A function that accepts a variable number of arguments. Defined with dotted parameter syntax: \texttt{(define (f x . rest) ...)}.

\item[vector]\index{vector}
A fixed-size, mutable, indexed collection with O(1) access by position.

\item[yield]\index{yield}
Voluntarily suspends the current fiber, allowing other fibers to run. Called implicitly during channel operations.

\end{description}
